% CSEC 620 Graduate Research Paper — Track B (reproduction / extension)
% Build on Overleaf (IEEEtran is built in) or locally: latexmk -pdf main.tex
% Tables/figures under generated/ and figures/ come from `make reproduce`
% (src/tables.py). Do not type result numbers by hand: use the tables or the
% macros in generated/numbers.tex (e.g. \EOneGistKnnAcc).
%
% Page budget (Milestone 2): 9 pages incl. ~1 page of references; cap is 10.
\documentclass[conference]{IEEEtran}
\usepackage{cite}
\usepackage{amsmath,amssymb}
\usepackage{graphicx}
\usepackage{booktabs}
\usepackage{multirow}
\usepackage{xcolor}
\usepackage{url}
\newcommand{\todo}[1]{\textcolor{red}{[TODO: #1]}}
\newcommand{\inputgen}[1]{\IfFileExists{generated/#1}{\input{generated/#1}}{\todo{run make reproduce to generate #1}}}
\newcommand{\figgen}[2]{\IfFileExists{figures/#1}{\includegraphics[width=#2]{figures/#1}}{\todo{missing figures/#1}}}
\IfFileExists{generated/numbers.tex}{% Generated by src/tables.py from results/metrics.json; do not edit.
\newcommand{\EOneGistKnnAcc}{97.46\%}
\newcommand{\EOneGistKnnFone}{93.56\%}
\newcommand{\EOneCnnAcc}{99.29\%}
\newcommand{\EOneCnnFone}{98.15\%}
\newcommand{\EOneByteHistLrAcc}{96.82\%}
\newcommand{\EOneByteHistLrFone}{97.29\%}
\newcommand{\ETwoGistKnnAcc}{94.95\%}
\newcommand{\ETwoGistKnnFone}{86.65\%}
\newcommand{\ETwoGistKnnAccAll}{84.29\%}
\newcommand{\ETwoGistKnnFoneAll}{73.63\%}
\newcommand{\ETwoCnnAcc}{99.03\%}
\newcommand{\ETwoCnnFone}{97.51\%}
\newcommand{\ETwoCnnAccAll}{87.91\%}
\newcommand{\ETwoCnnFoneAll}{82.72\%}
\newcommand{\ETwoByteHistLrAcc}{95.67\%}
\newcommand{\ETwoByteHistLrFone}{95.67\%}
\newcommand{\ETwoByteHistLrAccAll}{84.93\%}
\newcommand{\ETwoByteHistLrFoneAll}{80.67\%}
\newcommand{\EThreeGistKnnAcc}{90.43\%}
\newcommand{\EThreeGistKnnFone}{88.02\%}
\newcommand{\EThreeCnnAcc}{94.04\%}
\newcommand{\EThreeCnnFone}{93.77\%}
\newcommand{\EThreeByteHistLrAcc}{87.06\%}
\newcommand{\EThreeByteHistLrFone}{84.92\%}
}{}

\begin{document}

\title{\todo{Title}}
\author{\IEEEauthorblockN{\todo{Name}}
\IEEEauthorblockA{Department of Cybersecurity, Rochester Institute of Technology\\
Rochester, NY, USA\\ \todo{email}}}
\maketitle

% ~0.2 pages: problem; what was reproduced; the three extension findings WITH
% numbers (from generated/numbers.tex); one-line takeaway.
\begin{abstract}
\todo{Abstract}
\end{abstract}

\begin{IEEEkeywords}
malware classification, byteplot, GIST, reproducibility, data leakage, adversarial robustness
\end{IEEEkeywords}

% ---------------------------------------------------------------- I (0.8 p)
\section{Introduction}
% - Why static, execution-free classification matters.
% - The headline-accuracy puzzle: 98--99% everywhere [1], [4], [7], [8].
% - RQ1--RQ4 (below), contributions list, how this differs from [14].
\todo{Introduction}

\begin{itemize}
  \item \textbf{RQ1 (Reproduction):} Does GIST + $k$-NN reach roughly 98\% on Malimg under the protocol of~\cite{nataraj2011images}, and how close does a fine-tuned CNN come to~\cite{kalash2018mcnn,vasan2020imcfn} under the same split?
  \item \textbf{RQ2 (Leakage):} How much accuracy survives when exact and near-duplicate images are kept on one side of the split?
  \item \textbf{RQ3 (Generalization):} Do the method rankings hold on Microsoft BIG~2015 under an identical pipeline?
  \item \textbf{RQ4 (Robustness):} How quickly do GIST + $k$-NN and the CNN degrade under functionality-preserving byte appending?
\end{itemize}

% ---------------------------------------------------------------- II (0.6 p)
\section{Background}
% Byteplot construction + width table [1]; GIST [3]; CNN transfer learning [9];
% PE format basics needed for the threat model [12].
\todo{Background}

% ---------------------------------------------------------------- III (1.3 p)
\section{Related Work}
% Organise by pipeline stage following [17], NOT paper-by-paper:
% (a) image construction [1], [6]; (b) features/classifiers [4], [5], [7], [8], [10];
% (c) evaluation practice [15], [16]; (d) robustness [11]--[14].
\todo{Related work}

% Fair-comparison table (assignment requirement). Fill each cell from the paper
% itself. A row that differs from ours in ANY column is "indicative only".
% - Sample counts: give each published number's own dataset count (Malimg is
%   9,458 in [1] but 9,339 in our copy and in [14]).
% - Every BIG 2015 row is indicative only: our 2,042-sample capped subset is
%   more balanced than the full 10,868-sample training set used in [4], [5].
\begin{table*}[t]
\centering
\caption{Published results and whether each is comparable to ours. ``Indicative'' = differs from our setup in at least one column.}
\label{tab:published}
\resizebox{\textwidth}{!}{%
\begin{tabular}{lllllllll}
\toprule
Work & Dataset & Method & Samples & Split protocol & Preprocessing / size & Metric & Reported & Comparable? \\
\midrule
\cite{nataraj2011images} & Malimg & GIST + $k$-NN & 9,458 & \todo{} & \todo{} & Acc. & $\approx$98\% & Indicative (count) \\
\cite{kalash2018mcnn} & Malimg & VGG16 (M-CNN) & \todo{} & \todo{} & \todo{} & Acc. & 98.52\% & \todo{} \\
\cite{gibert2019cnn} & Malimg & CNN & \todo{} & \todo{} & \todo{resizing method} & \todo{} & \todo{} & \todo{} \\
\cite{vasan2020imcfn} & Malimg & Fine-tuned CNN, color & \todo{} & \todo{} & \todo{augmentation} & Acc. & 98.82\% & \todo{} \\
\cite{tariq2026framework} & Malimg & \todo{} & 9,339 & \todo{} & oversampling & \todo{} & \todo{} & \todo{} \\
\midrule
\cite{kalash2018mcnn} & BIG 2015 & VGG16 (M-CNN) & \todo{} & \todo{} & \todo{} & Acc. & 99.97\% & Indicative (subset) \\
\cite{gibert2019cnn} & BIG 2015 & CNN & \todo{} & \todo{} & \todo{} & \todo{} & \todo{} & Indicative (subset) \\
\midrule
This work (E1) & Malimg & see Table~\ref{tab:main} & 9,339 & stratified 10-fold & \todo{} & Acc., macro-F1 & Table~\ref{tab:main} & -- \\
This work (E3) & BIG 2015 & see Table~\ref{tab:main} & 2,042 & grouped 10-fold & width table [1] & Acc., macro-F1 & Table~\ref{tab:main} & -- \\
\bottomrule
\end{tabular}}
\end{table*}

% ---------------------------------------------------------------- IV (1.6 p)
\section{Methodology}
% Datasets and counts (report the Malimg version, count and SHA-256 you used);
% dedup procedure (pixel SHA-256 + near-duplicate threshold and how it was tuned);
% split protocols; GIST + k-NN details incl. every assumption [1], [3];
% CNN training [9]; threat model and attacks [11], [12]; metrics; make reproduce.
\subsection{Datasets}
% MUST INCLUDE (write in your own words):
% Malimg [1]
%  - Source: the exact mirror used (Kaggle page + version); the original UCSB
%    download link no longer works.
%  - Size: 9,339 images, 25 families. Differs from the 9,458 reported in [1];
%    recent work [14] also uses 9,339.
%  - Format: ships as pre-rendered grayscale PNGs, so the width table was NOT
%    applied by us for this dataset.
% Microsoft BIG 2015 \cite{ronen2018big2015}
%  - Source: the Kaggle competition; cite Ronen et al.
%  - Used: only the labelled training set (test labels were never released),
%    and only the .bytes hex dumps, not .asm. PE headers were removed by the
%    dataset's creators.
%  - "??" entries: DROPPED (removed before rendering, not mapped to 0); see
%    src/data.py: parse_bytes_file. Say this changes byte offsets slightly.
%  - Rendering: .bytes -> byteplot with [1]'s width table; last partial row
%    zero-padded (assumption); src/data.py: byteplot.
%  - Subset: up to 250 per family, all 42 Simda kept -> 2,042 samples.
%    State the rule: \todo{fixed seed N / method}. The id list is committed as
%    data/big2015_subset.txt and read by the loader (config data.big2015_ids),
%    so the exact subset can be recreated.
\begin{table}[t]
\centering
\caption{Datasets as used in this work.}
\label{tab:datasets}
\resizebox{\columnwidth}{!}{%
\begin{tabular}{lll}
\toprule
 & Malimg \cite{nataraj2011images} & BIG 2015 \cite{ronen2018big2015} \\
\midrule
Source & \todo{Kaggle mirror, version} & Kaggle competition (train set) \\
Samples used & 9,339 (9,458 in \cite{nataraj2011images}) & 2,042 of 10,868 \\
Families & 25 & 9 \\
Input format & pre-rendered grayscale PNG & \texttt{.bytes} hex dumps (no PE headers) \\
Rendering & none (as shipped) & width table of \cite{nataraj2011images}; \texttt{??} dropped \\
Sampling & all & $\le$250 per family; all 42 Simda \\
Used in & E1, E2, E4 & E3 \\
\bottomrule
\end{tabular}}
\end{table}
\todo{Datasets prose}
\subsection{Deduplication and Splits}
\todo{}
\subsection{Models}
\todo{}
\subsection{Threat Model and Attacks}
\todo{}
\subsection{Metrics and Statistics}
\todo{}

% ---------------------------------------------------------------- V (2.0 p)
\section{Results}
\begin{table}[t]
\centering
\caption{Mean $\pm$ std over seeds (out-of-fold, pooled per seed).}
\label{tab:main}
\resizebox{\columnwidth}{!}{\inputgen{table_main.tex}}
\end{table}

\subsection{RQ1: Reproduction}
\todo{}
\subsection{RQ2: Duplicates and Leakage}
\begin{table}[t]
\centering
\caption{Per-family duplicate rate and recall, E1 (random) vs.\ E2 (grouped).}
\label{tab:perfamily}
\resizebox{\columnwidth}{!}{\inputgen{table_per_family.tex}}
\end{table}
\todo{}
\subsection{RQ3: Cross-Dataset}
% MUST INCLUDE: BIG 2015 numbers are INDICATIVE ONLY against [4], [5]. Our
% capped subset (<=250/family) is more balanced than the full training set, so
% accuracy here is not directly comparable to full-dataset results. Point to
% the "Comparable?" column of Table~\ref{tab:published}.
\todo{}
\subsection{RQ4: Robustness}
\begin{figure*}[t]
\centering
\figgen{e4_budget_curves.png}{\textwidth}
\caption{Accuracy vs.\ append budget (mean $\pm$ std over seeds).}
\label{fig:budget}
\end{figure*}
\begin{table}[t]
\centering
\caption{Robustness AUC (\%), area under accuracy-vs-budget; 100 = no degradation.}
\label{tab:robust}
\resizebox{\columnwidth}{!}{\inputgen{table_robustness.tex}}
\end{table}
\todo{}

% ---------------------------------------------------------------- VI (0.9 p)
\section{Discussion}
% Is Malimg still a useful benchmark? Why GIST and CNN converge on easy data;
% robustness gap and deployment; [16]'s pitfalls applied to the literature [15], [16].
\todo{Discussion}

% ---------------------------------------------------------------- VII (0.4 p)
\section{Limitations and Threats to Validity}
% No timestamps (no temporal split) [15]; Malimg ships images not binaries, so
% functionality of perturbed files cannot be verified; small family count; compute.
% MUST INCLUDE:
%  - Unverified mirror: we cannot confirm the Malimg mirror is identical to the
%    authors' original data (count already differs: 9,339 vs 9,458).
%  - Storage limits: disk space, not a methodological choice, is why a BIG 2015
%    subset (2,042 of 10,868) was used instead of the full set.
%  - IF E3 IS DROPPED: say so here, and move cross-dataset testing to Future Work.
\todo{Limitations}

% ---------------------------------------------------------------- VIII (0.3 p)
\section{Conclusion and Future Work}
% Answers to RQ1--RQ4; stronger attacks [13], [14]; time-aware datasets.
% IF E3 IS DROPPED: cross-dataset testing on BIG 2015 becomes future work.
\todo{Conclusion}

\bibliographystyle{IEEEtran}
\bibliography{refs}

\appendices
\section{Code and Reproducibility}
% Repo URL + tag; `make reproduce` (config.yaml) vs `make quick`; which produced the numbers.
\todo{}
\section{AI Use and Statement of Original Work}
% Required by the assignment. Log every substantive AI use (see SUBMISSION.md).
\todo{}

\end{document}
